% Grado 4 -- generado por tools/build_tex.py a partir de raw/grado_04.txt
\section{Grado 4}\label{grado:4}

% PDF p. 30
\dba{1}{Interpreta las fracciones como razón, relación parte todo, cociente y operador en diferentes contextos.}

\evidencias

\begin{itemize}
  \item Describe situaciones en las cuales puede usar fracciones y decimales.
  \item Reconoce situaciones en las que dos cantidades covarían y cuantifica el efecto que los cambios en una de ellas tienen en los cambios de la otra y a partir de este comportamiento determina la razón entre ellas.
\end{itemize}

\ejemplo

Las limonadas de don Diego son famosas. Tienen un sabor característico, quien las haya probado es capaz de identificarlas en cualquier parte. Aunque no se conoce la receta de don Diego, si se sabe que él utiliza para un litro de agua, seis limones y tres cucharadas de azúcar.

\fig{g04_p30_1}{Receta con ingredientes y cantidades (tazas de agua, limones, cucharadas de azúcar) para escalar proporciones.}

\begin{itemize}
  \item Tratando de imitar la receta, Carlos preparó una limonada con dos jarras de agua (de litro cada una), diez limones y seis cucharadas de azúcar. Justifica si la limonada de Carlos tendría o no el mismo sabor que la de don Diego. En caso de considerar que no, menciona las diferencias que tendrían.
  \item Si don Diego dispone de 18 limones, averigua la cantidad de agua y azúcar que debería utilizar si quiere preparar su limonada.
  \item Propone otras posibilidades de preparar limonadas con el sabor característico de don Diego variando las cantidades de los ingredientes.
  \item Explica qué ocurre si las personas utilizan diferentes medidas para las cucharas de azúcar.
  \item Completa las instrucciones que deben darse para la preparación de otras limonadas que tengan la misma concentración que la de Don Diego.
\end{itemize}

% PDF p. 30
\dba{2}{Describe y justifica diferentes estrategias para representar, operar y hacer estimaciones con números naturales y números racionales (fraccionarios) 1 , expresados como fracción o como decimal}

\evidencias

\begin{itemize}
  \item Utiliza el sistema de numeración decimal para representar, comparar y operar con números mayores o iguales a 10.000.
  \item Describe y desarrolla estrategias para calcular sumas y restas basadas en descomposiciones aditivas y multiplicativas.
\end{itemize}
1 No se espera en este nivel escolar un estudio profundo de los números racionales como sistema numérico, sino una primera aproximación a las cantidades -generalmente llamada en los grados de primaria como fraccionarios- que expresan una razón entre otras dos cantidades, y cuyo resultado no siempre da un número exacto de unidades. Aunque estrictamente hablando los conceptos representados por estas expresiones no son coincidentes se acompañará la palabra “racionales” de la palabra “fraccionario” escrita entre paréntesis.

\begin{itemize}
  \item Utiliza y justifica algoritmos estandarizados y no estandarizados para realizar operaciones aditivas con representaciones decimales provenientes de fraccionarios cuyas expresiones tengan denominador 10, 100, etc.
  \item Identifica y construye fracciones equivalentes a una fracción dada.
  \item Propone estrategias para calcular sumas y restas de algunos fraccionarios.
\end{itemize}

\ejemplo

El banco de un juego de mesa tiene los siguientes billetes y fichas para comprar tarjetas de autos:

\fig{g04_p31_1}{Fichas con las fracciones 1/10, 1/2 y 1/100, y una ficha con el número 1.}

Cada billete representa una parte del valor de la ficha circular y ésta se le entrega únicamente a quien deposite \$10.000 en el banco del juego. Encuentra la cantidad de dinero que se puede canjear en el banco cuando tenga trece billetes azules, veintidós billetes rojos, dos billetes grises y seis fichas. Determina si con ese dinero es posible comprar la tarjeta de un auto cuyo precio en el juego es de 10 fichas circulares.


% PDF p. 31
\dba{3}{Establece relaciones mayor que, menor que, igual que y relaciones multiplicativas entre números racionales en sus formas de fracción o decimal.}

\evidencias

\begin{itemize}
  \item Construye y utiliza representaciones pictóricas para comparar números racionales (como fracción o decimales).
  \item Establece, justifica y utiliza criterios para comparar fracciones y decimales.
  \item Construye y compara expresiones numéricas que contienen decimales y fracciones.
\end{itemize}

\ejemplo

\fig{g04_p31_2}{Balanza en equilibrio con un cubo en un plato y esferas en el otro.}

un lado de la balanza se el otro lado un objeto con caras cuadradas. Según la información de la imagen:

\begin{itemize}
  \item Indica cuál de los dos objetos es más pesado.
  \item Describe la relación que hay entre los pesos de los dos objetos.
  \item Explica si las relaciones expresadas por las balanzas de la figura siguiente concuerdan con los datos de la imagen anterior. En caso que haya algunas que no sean correctas, dibuja la posición de la balanza que esté acorde con las cantidades presentadas.
\end{itemize}
\fig{g04_p32_1}{Balanza con un cubo y esferas.}

\fig{g04_p32_2}{Dos balanzas con distintas combinaciones de cubos y esferas para deducir equivalencias de peso.}


% PDF p. 32
\dba{4}{Caracteriza y compara atributos medibles de los objetos (densidad, dureza, viscosidad, masa, capacidad de los recipientes, temperatura) con respecto a procedimientos, instrumentos y unidades de medición; y con respecto a las necesidades a las que responden.}

\evidencias

\begin{itemize}
  \item Reconoce que para medir la capacidad y la masa se hacen comparaciones con la capacidad de recipientes de diferentes tamaños y con paquetes de diferentes masas, respectivamente (litros, centilitros galón, botella, etc., para capacidad, gramos, kilogramos, libras, arrobas, etc., para masa.)
  \item Diferencia los atributos medibles como capacidad, masa, volumen, entre otros, a partir de los procedimientos e instrumentos empleados para medirlos y los usos de cada uno en la solución de problemas.
  \item Identifica unidades y los instrumentos para medir masa y capacidad, y establece relaciones entre ellos.
  \item Describe procesos para medir capacidades de un recipiente o el peso de un objeto o producto.
  \item Argumenta sobre la importancia y necesidad de medir algunas magnitudes como densidad, dureza, viscosidad, masa, capacidad, etc.
\end{itemize}

\ejemplo

En clase de sociales le enseñan a Felipe que es conveniente seleccionar productos que, además de ser económicos, ofrezcan posibilidades de reciclaje, por el tipo de material del empaque.

¿Qué criterios son adecuados para seleccionar entre varias marcas el mejor producto por economía y posibilidades de reciclaje?

Compara la información brindada en los empaques de dos o más productos para tomar decisiones, cuando la información no es suficiente propone procedimientos de medida para hacer las comparaciones.


% PDF p. 33
\dba{5}{Elige instrumentos y unidades estandarizadas y no estandarizadas para estimar y medir longitud, área, volumen, capacidad, peso y masa, duración, rapidez, temperatura, y a partir de ellos hace los cálculos necesarios para resolver problemas.}

\evidencias

\begin{itemize}
  \item Expresa una misma medida en diferentes unidades, establece equivalencias entre ellas y toma decisiones de la unidad más conveniente según las necesidades de la situación.
  \item Propone diferentes procedimientos para realizar cálculos (suma y resta de medidas, multiplicación y división de una medida y un número) que aparecen al resolver problemas en diferentes contextos.
  \item Emplea las relaciones de proporcionalidad directa e inversa para resolver diversas situaciones.
  \item Propone y explica procedimientos para lograr mayor precisión en la medición de cantidades de líquidos, masa, etc.
\end{itemize}

\ejemplo

La receta de la torta de vainilla para 20 personas es 1

\fig{g04_p33_1}{Receta: ingredientes (1 taza de azúcar, 125 g de mantequilla, 2 huevos, 2 cucharadas de esencia de vainilla, 1½ tazas de harina, 1½ cucharaditas de polvo de hornear) y pasos de preparación.}

El azúcar, la harina y la mantequilla se venden por libras. Identifica qué cantidad de azúcar y qué cantidad de mantequilla (en libras) se requiere para hacer la torta. Propone formas más precisas para medir las cantidades de leche, de esencia de vainilla y de otros ingredientes; y establece las cantidades. Determina los grados Fahrenheit a los que se debe programar el horno para hornear la torta y las cantidades de cada ingrediente que se requieren para elaborar la receta con las mismas características de sabor para 30 personas.

1 Tomado de http://allrecipes.com.mx/receta/157/pastel-simplemente-blanco.aspx


% PDF p. 33
\dba{6}{Identifica, describe y representa figuras bidimensionales y tridimensionales, y establece relaciones entre ellas.}

\evidencias

\begin{itemize}
  \item Arma, desarma y crea formas bidimensionales y tridimensionales.
  \item Reconoce entre un conjunto de desarrollos planos, los que corresponden a determinados sólidos atendiendo a las relaciones entre la posición de las diferentes caras y aristas.
\end{itemize}

\ejemplo

Construye esculturas geométricas con cubos y prismas triangulares (medios cubos) y representa de manera bidimensional la representación tridimensional.

\fig{g04_p34_1}{Figuras (animales) construidas con cubos de colores.}

\fig{g04_p34_2}{Figura formada por cuadrados sobre cuadrícula (desarrollo plano / área por conteo).}


% PDF p. 34
\dba{7}{Identifica los movimientos realizados a una figura en el plano respecto a una posición o eje (rotación, traslación y simetría) y las modificaciones que pueden sufrir las formas (ampliación- reducción).}

\evidencias

\begin{itemize}
  \item Aplica movimientos a figuras en el plano.
  \item Diferencia los efectos de la ampliación y la reducción.
  \item Elabora argumentos referente a las modificaciones que sufre una imagen al ampliarla o reducirla.
  \item Representa elementos del entorno que sufren modificaciones en su forma.
\end{itemize}

\ejemplo

La familia de Francisco estuvo de vacaciones en la finca de los abuelos. Para guardar un recuerdo tomaron una fotografía del lugar.

\fig{g04_p34_3}{Fotografía de un paisaje de granja en un marco de película, tamaño original.}

Francisco quiere poner la fotografía en su habitación para recordar sus vacaciones, pero debe disminuir el tamaño de la imagen. Escoge la imagen que representa una reducción de la foto, justifica y describe el procedimiento realizado para seleccionar la imagen.

\fig{g04_p34_5}{Versión (a) de la fotografía del paisaje.}

\fig{g04_p34_4}{Versiones (b) y (d) de la fotografía del paisaje con otras dimensiones (ampliación/reducción, deformación).}

\fig{g04_p34_6}{Versión (c) de la fotografía del paisaje, ancha y baja.}

(c)

Dibuja la finca del abuelo pero dos veces más grande que la que aparece en la fotografía.


% PDF p. 35
\dba{8}{Identifica, documenta e interpreta variaciones de dependencia entre cantidades en diferentes fenómenos (en las matemáticas y en otras ciencias) y los representa por medio de gráficas.}

\evidencias

\begin{itemize}
  \item Realiza cálculos numéricos, organiza la información en tablas, elabora representaciones gráficas y las interpreta.
  \item Propone patrones de comportamiento numérico.
  \item Trabaja sobre números desconocidos y con esos números para dar respuestas a los problemas.
\end{itemize}

\ejemplo

Consigue una pelota pequeña y mide la altura ‘h’ (medida en centímetros) hasta la que rebota cuando se deja caer sobre una superficie dura (cemento) desde diversas alturas H, medida en centímetros (Figura 1). Realiza un registro aproximado de la altura H desde la que se suelta la pelota, así como de la altura h a la que rebota. Representa de otras formas la relación que encuentra entre la altura H (inicial) y la altura h (alcanzada en cada rebote). Escribe algunas conclusiones de esta exploración.

\fig{g04_p35_1}{Figura 1: poste vertical graduado con una pelota que cae, marcando alturas.}

\fig{g04_p35_2}{Trayectoria de una pelota que rebota, con posiciones sucesivas.}


% PDF p. 35
\dba{9}{Identifica patrones en secuencias (aditivas o multiplicativas) y los utiliza para establecer generalizaciones aritméticas o algebraicas.}

\evidencias

\begin{itemize}
  \item Comunica en forma verbal y pictórica las regularidades observadas en una secuencia.
  \item Establece diferentes estrategias para calcular los siguientes elementos en una secuencia.
  \item Conjetura y argumenta un valor futuro en una secuencia aritmética o geométrica (por ejemplo, en una secuencia de figuras predecir la posición 10, 20 o 100)
\end{itemize}

\ejemplo

Explora el efecto que tiene el signo “=” (igual) sobre el resultado a medida que se presiona varias veces.

\fig{g04_p35_3}{Calculadora científica con la que se exploran el signo igual y las operaciones repetidas.}

a. Describe y compara el efecto que la acción descrita tiene, si se presiona varias veces el signo igual después de digitar el símbolo de la multiplicación o el símbolo de la división. b. Describe las operaciones y resultados que muestra la calculadora, cuando se presiona 4 x 2 = = = = = =. c. Se digita la operación 3 X 4 y luego se presiona la tecla igual diez veces ¿será posible obtener un número menor que 1.000? Estima un número aproximado de veces que deberías presionar el “igual” para obtener el resultado más cercano a 1.000. Utiliza la calculadora para realizar las operaciones y verificar el resultado. d. Determina el mínimo número de veces que se debe presionar el signo igual después de hacer la operación 2048÷2 para obtener un número no natural.


% PDF p. 36
\dba{10}{Recopila y organiza datos en tablas de doble entrada y los representa en gráficos de barras agrupadas o gráficos de líneas, para dar respuesta a una pregunta planteada. Interpreta la información y comunica sus conclusiones.}

\evidencias

\begin{itemize}
  \item Elabora encuestas sencillas para obtener la información pertinente para responder la pregunta.
  \item Construye tablas de doble entrada y gráficos de barras agrupadas, gráficos de líneas o pictogramas con escala.
  \item Lee e interpreta los datos representados en tablas de doble entrada, gráficos de barras agrupados, gráficos de línea o pictogramas con escala.
  \item Encuentra e interpreta la moda y el rango del conjunto de datos y describe el comportamiento de los datos para responder las preguntas planteadas.
\end{itemize}

\ejemplo

La siguiente información fue recolectada en un hato lechero. Con dicha información elabora un informe, para enviarlo al dueño del hato, en el que se compara la producción de leche en horas de la mañana y en horas de la tarde, así como la variación de la producción por vaca.

\fig{g04_p36_1}{Vaca y tabla de datos por categorías (producción registrada).}


% PDF p. 36
\dba{11}{Comprende y explica, usando vocabulario adecuado, la diferencia entre una situación aleatoria y una determinística y predice, en una situación de la vida cotidiana, la presencia o no del azar}

\evidencias

\begin{itemize}
  \item Reconoce situaciones aleatorias en contextos cotidianos.
  \item Enuncia diferencias entre situaciones aleatorias y deterministas.
  \item Usa adecuadamente expresiones como azar o posibilidad, aleatoriedad, determinístico.
  \item Anticipa los posibles resultados de una situación aleatoria.
\end{itemize}

\ejemplo

En las siguientes situaciones reconoce la presencia o no del azar y expone diferencias entre ellas para expresar la posibilidad de conocer, con exactitud, los resultados que se tendrán antes de la ocurrencia del evento. a. La selección de la cancha que le corresponde a uno de los equipos cuando se inicia un partido de fútbol en el campeonato mundial. b. La selección del nombre del mes entrante. c. La selección del menú del refrigerio de la mañana. d. La conformación de dos equipos para jugar fútbol.


